\documentclass[11pt]{article}
% Standalone preamble for web-only supplementary notes.
\usepackage[T1]{fontenc}
\usepackage{amsmath,amssymb,amsthm}
\usepackage{hyperref}
\usepackage{cleveref}
\DeclareMathOperator{\Map}{Map}
\DeclareMathOperator{\Fun}{Fun}
\newcommand{\iso}{\xrightarrow{\sim}}
\newtheorem{lemma}{Supplementary Lemma}
\newtheorem{corollary}[lemma]{Supplementary Corollary}
\crefname{lemma}{Supplementary Lemma}{Supplementary Lemmas}
\Crefname{lemma}{Supplementary Lemma}{Supplementary Lemmas}
\crefname{corollary}{Supplementary Corollary}{Supplementary Corollaries}
\Crefname{corollary}{Supplementary Corollary}{Supplementary Corollaries}

\begin{document}
\section*{The attaching maps of a join}

The \href{joins.html}{mapping-out construction of a join} glues a cylinder
to its two endpoint categories. To construct the map induced by functors
on those categories, the two attaching squares must be compared with their
specified restrictions. We describe this calculation for absolute joins.
Assume the Chapter 1 structures, including universal coproducts, the walking
morphism and the mapping-out join axiom.

For $C,D$, put $W=C\times D$ and write
\[
 B=W\sqcup W,\qquad Z=W\times[1],\qquad E=C\sqcup D.
\]
Let $\iota_e=(\operatorname{id}_W,e)\colon W\to Z$ be endpoint insertion.
The two attaching maps are
\[
 t=[\iota_0,\iota_1]\colon B\to Z,\qquad
 b=[\operatorname{in}_C\operatorname{pr}_C,
       \operatorname{in}_D\operatorname{pr}_D]\colon B\to E.
\]

\begin{lemma}
\label[lemma]{lem:Supplementary_Join_Product_Coordinates}
The chosen join $C\star D$ is the pushout of $t$ and $b$.
Its boundary map $a\colon E\to C\star D$ is the original join inclusion;
its cylinder map $j\colon Z\to C\star D$ carries a specified matching
$jt=ab$.
\end{lemma}
\begin{proof}[Proof idea]
The comparison $C\times D\to C\times_\star D$ is an equivalence with
specified projection computations. Use it on the cylinder, and use
distributivity to replace its boundary by two copies of $W$.
The projection computations identify the two attaching maps with $t,b$.
Transport the whole universal cocone through these equivalences.
Finally identify its boundary leg with the original join inclusion and
transport the matching along that identification. This gives the claimed
pushout with its specified legs.
\end{proof}

\begin{lemma}
\label[lemma]{lem:Supplementary_Join_Attaching_Comparisons}
For $f\colon C\to C'$ and $g\colon D\to D'$, put
$v=f\times g$, $d=v\sqcup v$, $z=v\times\operatorname{id}_{[1]}$ and
$e=f\sqcup g$. There are comparisons
\[
 t'd=zt,\qquad b'd=eb
\]
whose restrictions to the two summands are the endpoint-insertion and
projection comparisons.
\end{lemma}
\begin{proof}[Proof idea]
On each summand the first equation reduces to
$\iota'_{\varepsilon}v=z\iota_{\varepsilon}$ for $\varepsilon=0,1$. Compare the product projections: the $W'$ projection
is $v$ and the interval projection is the same constant endpoint.
For the second equation, compute the $C'$ and $D'$ projections of $f\times g$
and then apply the appropriate coproduct inclusion.

These calculations initially compare the restrictions of $t'd,zt$ and
$b'd,eb$. Lift each pair through the equivalence on identification animae
induced by coproduct restriction. Its two computation laws retain the
specified comparisons on each summand. This is the
\href{supplementary-identifications-on-coproducts.html}{prescribed-image
coproduct calculation}.
\end{proof}

\begin{lemma}
\label[lemma]{lem:Supplementary_Join_Cocone_Computation}
There is a functor $F\colon C\star D\to C'\star D'$ with comparisons
\[
 Fj=j'z,\qquad Fa=a'e,
\]
compatible with the specified matchings on $B$.
\end{lemma}
\begin{proof}[Proof idea]
Restrict the target join cocone along the three maps $d,z,e$, using the
two comparisons of \Cref{lem:Supplementary_Join_Attaching_Comparisons}.
This is a cocone on the source span. Extend it through the source pushout.
The computation law of the extension compares the entire resulting
cocone with the one just constructed. Its two components are the displayed
identifications, and its matching equation is their required compatibility.
\end{proof}

This is the pushout-defined action on joins. The checked declarations linked
here do not supply its identity and composition laws or identify it with
the separate action constructed using dependent products. General
associativity of joins is also outside the scope of this note.
\end{document}
